\section{Opérateur champ}

    Cette section a pour objectif de généraliser l’usage des espaces de Fock aux cas d’espaces de bases continues.

    \subsection{Définition de l’opérateur champ}

    L’\textbf{opérateur champ} est défini comme opérateur d’annihilation $\hat{a}_{k}$, mais relatif à une base d’états individuels faisant intervenir l’indice de position continu $\mathbf{r}$ plutôt qu’un indice discret $k$. La relation de changement de base discrète (d’une base $\left\{\ket{u_k}\right\}$ à $\left\{v_{\ell}\right\}$) est :
    \[ \hat{a}_{v_{\ell}} = \sum_{k} \spr{v_{\ell}}{u_k} \hat{a}_{u_k} \]   

        \subsubsection{Définition}

        On commence par définir l’opérateur champ pour des particules sans spin avant de généraliser.

            \paragraph{Particules sans spin}

            Considérons la base des positions $\left\{\ket{\mathbf{r}}\right\}$. L’opérateur $\hat{a}_{\ell}$ devient alors un opérateur dépendant d’un indice continu $\mathbf{r}$. La notation que nous utiliserons est $\Psi(\mathbf{r})$, que nous définissons par analogie au cas discret, en posant $u_k(\mathbf{r}) = \spr{\mathbf{r}}{u_k}$ :
            \[ \Psi(\mathbf{r}) = \sum_{k} u_k(\mathbf{r}) \hat{a}_{k} \]
            De la même façon que l’opérateur $\hat{a}_{\ell}$ annihile une particule dans l’état $\ket{v_{\ell}}$, l’opérateur $\Psi(\mathbf{r})$ annihile une particule au point $\mathbf{r}$. 

            Son conjugué hermitique est l’opérateur de création d’une particule au point $\mathbf{r}$, et a pour expression 
            \[ \Psi^{\dagger}(\mathbf{r}) = \sum_{k} u_k^*(\mathbf{r}) \hat{a}_k^{\dagger} \]   
            On peut vérifier son action sur le ket vide :
            \begin{align*}
                \Psi^{\dagger}(\mathbf{r})\ket{0} 
                &=  \sum_{k} u_k^*(\mathbf{r}) \hat{a}_k^{\dagger} \ket{0} \\
                &=  \sum_{k} u_k^*(\mathbf{r}) \ket{u_k} \\
                &=  \sum_{k} \ket{u_k} \spr{u_k}{\mathbf{r}} \\
                &= \ket{\mathbf{r}}
            \end{align*}

            On peut inverser les formules pour obtenir l’expression des opérateurs initiaux :
            \begin{align*}
                \int d^3 r u_k^*(\mathbf{r}) \Psi(\mathbf{r}) 
                &= \int d^3 r u_k^*\mathbf{r} \sum_{k'} u_{k'}(\mathbf{r}) \hat{a}_{k'} \\
                &= \sum_{k'} \delta_{kk'} \hat{a}_{k'} \\
                &= \hat{a}_{k} \\
                \int d^3 r u_k(\mathbf{r}) \Psi^{\dagger}(\mathbf{r}) 
                &= \hat{a}_k^{\dagger}
            \end{align*}

            \paragraph{Particules à spin}

            On se place désormais dans la base $\left\{\ket{\mathbf{r}, \nu}\right\}$ où $\nu$ est l’inde de spin qui peut prendre $2 S + 1$ valeurs discrètes distinctes entre $-S$ et $S$. On doit donc maintenant ajouter une sommation sur les valeurs de l’indive $\nu$ de spin :
            \[ \ket{u_i} = \sum_{\nu = -S}^S \int d^3r u_k(\mathbf{r}, \nu) \ket{\mathbf{r}, \nu} \]
            où $u_k(\mathbf{r}, \nu) = u_k^{\nu}(\mathbf{r}) = \spr{\mathbf{r}, \nu}{u_k}$. La seconde écriture permet de se rappeler que $\mathbf{r}$ et $\nu$ ne jouent pas le même rôle car l’un est continu et l’autre discret. De façon inverse, on a 
            \[ \Psi_{\nu}(\mathbf{r}) = \sum_{k} u_k^{\nu}(\mathbf{r}) \hat{a}_{k} \]   
            Cet opérateur champ à composante de spin détruit une particule de spin $\nu$ au point $\mathbf{r}$. Son conjugué hermitique créé une telle particule, avec l’expression 
            \[ \Psi_{\nu}^{\dagger}(\mathbf{r}) = \sum_{k} u_k^{\nu *}(\mathbf{r}) \hat{a}^{\dagger}_{k} \]

        \subsubsection{Relations de commutation ou d’anticommutation}

        Ces relations sont analogues à celles rencontrés dans le cas discret. On pose $\left[\hat{A}, \hat{B}\right]_{-\eta} = AB - \eta BA$ où $\eta = 1$ pour des bosons et $\eta = -1$ pour des fermions.

            \paragraph{Particules sans spin}

            \begin{align*}
                \left[\Psi(\mathbf{r}), \Psi(\mathbf{r}')\right]_{- \eta} 
                &= \sum_{i,j} u_i(\mathbf{r}) u_j(\mathbf{r}') \left[\hat{a_i}, \hat{a_j}\right]_{- \eta} = 0 \\
                \left[\Psi^{\dagger}(\mathbf{r}), \Psi^{\dagger}(\mathbf{r}')\right]_{- \eta} 
                &= \sum_{i,j} u_i^*(\mathbf{r}) u_j^*(\mathbf{r}') \left[\hat{a_i}^{\dagger}, \hat{a_j}^{\dagger}\right]_{- \eta} = 0 \\ 
                \left[\Psi(\mathbf{r}), \Psi^{\dagger}(\mathbf{r}')\right]_{- \eta} 
                &= \sum_{i,j} u_i(\mathbf{r}) u_j^*(\mathbf{r}') \left[\hat{a_i}, \hat{a_j}^{\dagger}\right]_{- \eta} \\ 
                &= \sum_{i,j} u_i(\mathbf{r}) u_j^*(\mathbf{r}') \\
                &= \spr{\mathbf{r}}{\mathbf{r}'} = \delta(\mathbf{r} - \mathbf{r}') \\
            \end{align*}

            \paragraph{Particules avec spin}

            Les deux premières égalités restent vraies, et seulement la dernière impose aussi que les spins soient égaux :

            \begin{align*}
                \left[\Psi_{\nu}(\mathbf{r}), \Psi_{\nu'}(\mathbf{r}')\right]_{- \eta} 
                &= \sum_{i,j} u_i^{\nu}(\mathbf{r}) u_j^{\nu'}(\mathbf{r}') \left[\hat{a_i}, \hat{a_j}\right]_{- \eta} = 0 \\
                \left[\Psi_{\nu}^{\dagger}(\mathbf{r}), \Psi_{\nu'}^{\dagger}(\mathbf{r}')\right]_{- \eta} 
                &= \sum_{i,j} u_i^{\nu *}(\mathbf{r}) u_j^{\nu'*}(\mathbf{r}') \left[\hat{a_i}^{\dagger}, \hat{a_j}^{\dagger}\right]_{- \eta} = 0 \\ 
                \left[\Psi_{\nu}(\mathbf{r}), \Psi_{\nu'}^{\dagger}(\mathbf{r}')\right]_{- \eta} 
                &= \sum_{i,j} u_i^{\nu}(\mathbf{r}) u_j^{\nu'*}(\mathbf{r}') \left[\hat{a_i}, \hat{a_j}^{\dagger}\right]_{- \eta} \\ 
                &= \sum_{i,j} u_i^{\nu}(\mathbf{r}) u_j^{\nu'*}(\mathbf{r}') \\
                &= \spr{\mathbf{r}, \nu}{\mathbf{r}', \nu'} = \delta_{\nu \nu'} \delta(\mathbf{r} - \mathbf{r}') \\
            \end{align*}

    \subsection{Opérateurs symétriques}

    Dans la section précédente, nous avons exprimé les opérateurs symétriques à une particule à l’aide des opérateurs création et annihilation. Dans une base continue, on utilise alors l’opérateur champ et son conjugué hermitique.


       
